\begin{helper}
Fix one blue support label \(B\), one red label \(J\), a terminal
coordinate set \(U\), and a nodup terminal order \(o\) with
\(\operatorname{toFinset}(o)=U\).  Let \(\mathcal T\) be the finite
branch-label type and let \(\mathcal P\) be the finite transcript family.
For a realized pair \(c=(\beta,\tau)\), write
\[
R_c=R_\beta,\qquad P_c=P_\tau,\qquad A_c=A_\tau,\qquad Z_c=Z_\tau .
\]
Assume \(0\le p\le1/2\), \(|B|=u_B\), and \(B\subseteq U\).  Assume that
every realized pair with \(\tau\in\mathcal P\) satisfies the fixed
\((B,J)\) geometry:
\[
|R_\beta|=u_R,\qquad R_\beta\subseteq U,\qquad B\cap R_\beta=\emptyset,
\]
\[
P_\tau,A_\tau\subseteq U,\qquad P_\tau\cap A_\tau=\emptyset,\qquad
B\cup R_\beta\subseteq P_\tau,
\]
that \(P_\tau,A_\tau\) encode the complete terminal blue trace of \(\beta\),
and that
\[
Z_\tau\subseteq A_\tau,\qquad
\max(0,a-u_R-u_B)\le |Z_\tau|.
\]

Assume in addition the lower-level reveal-order residual certificate for
this fixed pair \(B,J\).  Namely, there is a finite residual leaf type
\(\Lambda\), a residual leaf assignment
\[
\iota:\mathcal T\times\mathcal P\to\Lambda,
\]
and a leaf mass \(\nu:\Lambda\to\mathbb R\).  The leaf masses are
nonnegative and normalized:
\[
0\le \nu(\lambda)\quad(\lambda\in\Lambda),\qquad
\sum_{\lambda\in\Lambda}\nu(\lambda)\le1.
\]
The leaf assignment is injective on realized cells: if
\(\tau,\tau'\in\mathcal P\), both \((\beta,\tau)\) and
\((\beta',\tau')\) are realized, and
\[
\iota(\beta,\tau)=\iota(\beta',\tau'),
\]
then \(\beta=\beta'\) and \(\tau=\tau'\).  Finally, every realized pair has
the residual leaf domination
\[
p^{|P_\tau|-u_R-u_B}(1-p)^{|A_\tau|-|Z_\tau|}
  \le \nu(\iota(\beta,\tau)).
\]
Here \(\Lambda,\iota,\nu\) are leaf-level scan data: they are not the
cell-indexed residual masses being constructed below.

Then there are residual masses \(\sigma_{\beta,\tau}\), indexed by formal
branch/transcript pairs, such that realized allowed pairs have
\[
0\le\sigma_{\beta,\tau},
\]
the realized residual mass is a subprobability,
\[
\sum_{\beta}\sum_{\tau\in\mathcal P:\,(\beta,\tau)\text{ realized}}
  \sigma_{\beta,\tau}\le1,
\]
and every realized allowed pair satisfies
\[
p^{|P_\tau|-u_R-u_B}(1-p)^{|A_\tau|-|Z_\tau|}
  \le \sigma_{\beta,\tau}.
\]
\end{helper}
\begin{proof}
Define the cell residual mass by reading the mass of the residual leaf owned
by that cell:
\[
\sigma_{\beta,\tau}:=\nu(\iota(\beta,\tau)).
\]
This is a formal definition for every branch/transcript pair; no
realization condition is needed to define it.

For a realized pair \((\beta,\tau)\), nonnegativity follows immediately from
the lower-level certificate:
\[
0\le \nu(\iota(\beta,\tau))=\sigma_{\beta,\tau}.
\]
The termwise residual lower bound is exactly the leaf-domination clause of
the certificate, after substituting the definition of \(\sigma\):
\[
p^{|P_\tau|-u_R-u_B}(1-p)^{|A_\tau|-|Z_\tau|}
  \le \nu(\iota(\beta,\tau))
  = \sigma_{\beta,\tau}.
\]

It remains to prove the subprobability bound.  Let
\[
\mathcal C=\{(\beta,\tau):\tau\in\mathcal P,\,
(\beta,\tau)\text{ is realized}\}
\]
be the finite set of realized cells.  By the injectivity hypothesis, the map
\[
(\beta,\tau)\mapsto\iota(\beta,\tau)
\]
is injective on \(\mathcal C\).  Therefore summing \(\sigma_{\beta,\tau}\)
over realized cells is the same as summing \(\nu\) over the image
\(\iota(\mathcal C)\), with no residual leaf counted twice:
\[
\sum_{(\beta,\tau)\in\mathcal C}\sigma_{\beta,\tau}
=
\sum_{\lambda\in\iota(\mathcal C)}\nu(\lambda).
\]
The image \(\iota(\mathcal C)\) is a subset of \(\Lambda\).  Since every
leaf mass is nonnegative, enlarging the sum from this image to all residual
leaves can only increase it:
\[
\sum_{\lambda\in\iota(\mathcal C)}\nu(\lambda)
\le
\sum_{\lambda\in\Lambda}\nu(\lambda).
\]
The lower-level reveal-order product-tree certificate gives
\[
\sum_{\lambda\in\Lambda}\nu(\lambda)\le1.
\]
Combining these three displayed relations proves
\[
\sum_{\beta}\sum_{\tau\in\mathcal P:\,(\beta,\tau)\text{ realized}}
  \sigma_{\beta,\tau}\le1.
\]
Thus the leaf-level product-tree data yield the required cell-indexed
residual certificate.
\end{proof}
